\documentclass[../Thesis]{subfiles}
\begin{document}

\chapter{Conclusions} \label{chapter5}
\chaptermark{Conclusions}
\section{Summary}
This dissertation set out to investigate the challenge of identifying copyrighted audio embedded within the rapidly growing domain of short-form video content. Traditional audio-fingerprinting systems, though effective for clean or minimally distorted recordings, consistently fail when audio has been transformed by being pitch-shifted, time-stretched, cropped, layered or mixed. These manipulations are inherent to short-form media platforms such as TikTok, YouTube Shorts, and Instagram Reels, where creators frequently modify audio either intentionally or as part of the content-creation workflow.

WaveID was proposed as a robust, learning-based approach to this emerging problem. Through a review of classical fingerprinting, deep audio embeddings, and recent advances in contrastive and self-supervised learning, the work identified a gap in the academic landscape: no current open system explicitly focuses on recognition of heavily transformed audio. The methodology for WaveID addressed this gap by combining openly licensed datasets, controlled synthetic transformations, deep-learning-based embedding extraction, and FAISS-based similarity search.

The feasibility investigation confirmed that preliminary prototypes can successfully extract meaningful embeddings from both unaltered and transformed audio samples. Early evidence demonstrates that contrastive learning approaches, particularly those trained in realistic audio augmentations, offer significantly more robustness than traditional fingerprints or "off-the-shelf" embedding models. These findings support the viability of WaveID as a foundation for more advanced research.

\section{Research Contributions}
The main contributions of this dissertation are summarised as follows:

\begin{itemize}
    \item A detailed analysis of the shortcomings of classical audio fingerprinting with respect to short-form transformations.
    \item A structured evaluation of deep audio embeddings and their sensitivity to pitch, tempo, and segment-length variations.
    \item The design of a methodological framework for training contrastive audio-embedding models tailored to distorted audio.
    \item A feasibility-tested prototype demonstrating that embedding-based matching remains viable even under significant audio transformations.
    \item A reproducible experimental foundation suitable for further development into a full-scale audio-identification system.
\end{itemize}

\section{Future Work}
While the feasibility results are encouraging, several directions for future exploration remain:

\begin{itemize}
    \item \textbf{Large-scale training:} Training contrastive models on millions of audio segments, including user-generated short-form content.
    \item \textbf{Advanced architectures:} Integrating transformer-based audio encoders (e.g., AST, HTS-AT) that may outperform CNN-based embeddings.
    \item \textbf{Multi-modal integration:} Combining audio and visual cues for improved robustness in real-world short-form videos.
    \item \textbf{Real-platform testing:} Evaluating WaveID against real UGC samples taken from platforms such as TikTok or Instagram Reels.
    \item \textbf{Latency optimisation:} Ensuring fast inference for real-time or near-real-time rights-management systems.
\end{itemize}

\section{Final Remarks}
The rise of creator-centric media demands modern, transformation-tolerant copyright-identification techniques. WaveID demonstrates a promising first step in meeting this requirement, offering a reproducible framework built on contemporary deep-learning methods. While substantial work remains, this dissertation lays the technical and conceptual groundwork for next-generation copyright-protection systems capable of operating in the dynamic landscape of short-form video content.

\end{document}
